\documentclass[10pt,a4paper]{amsart}
\usepackage[margin=19mm]{geometry}
\usepackage{amsmath,amssymb,mathtools}
\usepackage[scr=rsfs,cal=euler]{mathalfa}
\usepackage{xfrac}
\usepackage[unicode,hidelinks,bookmarks=false]{hyperref}
\newcommand{\ie}{i.e.\ }
\newcommand{\cf}{cf.\ }
\newcommand{\resp}{respectively\ }
\newcommand{\Subsection}{\subsection}
\providecommand{\nobreakdash}{\nobreak\hbox{-}}
\setlength{\emergencystretch}{2em}
\multlinegap0pt
\usepackage{amssymb}
\usepackage{mathtools}
\usepackage{xcolor}
\usepackage[shortlabels]{enumitem}
\usepackage{tikz}
\newtheorem{lemma}{Lemma}
\newcommand{\Om}{\Omega}
\newcommand{\dive}{\operatorname{div}}
\title{The Quintic Wave Equation with Kelvin-Voigt Damping: Strichartz Estimates, Well-posedness and Global Stabilization}
\author{Marcelo Moreira Cavalcanti, Valeria Neves Domingos Cavalcanti}
\date{Source: arXiv:2603.02009v3 (CC BY 4.0)}
\begin{document}
\maketitle

\noindent\emph{Source and licence.} The Quintic Wave Equation with Kelvin-Voigt Damping: Strichartz Estimates, Well-posedness and Global Stabilization. Version arXiv:2603.02009v3 (CC BY 4.0), distributed by arXiv under the Creative Commons Attribution 4.0 International licence, http://creativecommons.org/licenses/by/4.0/, as recorded in the arXiv licence metadata for this exact version and shown on the abstract page https://arxiv.org/abs/2603.02009v3. Full licence notice: \texttt{licenses/arxiv-2603.02009-CC-BY-4.0.txt}.\par\smallskip

\noindent\textit{Editorial scope.} Counted unit: the article's lemma lem:chain (source0.tex lines 1533--1545). The article's complete original proof of it is delivered with it (source0.tex lines 1547--1606, one \texttt{proof} environment of 60 lines). No mathematics is written, repaired or re-derived: the statement and the proof are the article's own lines, in the article's own notation, and no retained line is substituted or re-flowed.  Retained uncounted as the source-local context the proof needs: lines 312--318; lines 1506--1510; lines 1516--1518.  Nothing was omitted from any retained span, so the package carries no omission record.  No retained line contains a citation, so the wrapper carries no bibliography and no bibliography entry.  No retained line contains a figure environment, an image-inclusion command or drawing code, so no image is delivered and the package declares no asset dependency.  Editorial formatting only, in addition: an AMS \texttt{amsart} preamble that restates the article's own declarations and macros used by the retained lines, namely the environments \texttt{lemma} on separate counters, together with the standard packages those lines need.  The retained environments print in this document as Lemma 1 in order of appearance, whereas the article prints the same items with the numbering of its own full document.  The complete native source of the article as distributed in the arXiv e-print for this exact version is bundled unchanged as \texttt{sources/195588/source0.tex}.
\begin{equation}\label{eq:main}
\begin{cases}
u_{tt} - \Delta u - \hbox{div}(a(x) \nabla u_t) + u^5 = 0 & \text{in } \Omega \times (0,T), \\
u = 0 & \text{on } \partial \Omega \times (0,T), \\
u(0) = u_0, \quad u_t(0) = u_1 & \text{in } \Omega,
\end{cases}
\end{equation}

\begin{equation}\label{eq:two-bounds}
\|u\|_{L^\infty(0,T^*;H_0^1)}\le\sqrt{2E_0},~
\|u_t\|_{L^\infty(0,T^*;L^2)}\le\sqrt{2E_0},~
\int_0^{T^*}\!\!\!\int_\Om a|\nabla u_t|^2\,dx\,dt\le E_0
\end{equation}

\begin{equation}\label{eq:astar}
a_*:=\min_{\overline{\omega_2}}a>0
\end{equation}

\begin{lemma}[Local parabolic time-regularity of the velocity]\label{lem:chain}
Under the above hypotheses, with $T:=T^*$ if $T^*<\infty$ and $T$ arbitrary finite otherwise:
\begin{enumerate}
\item[\rm(i)] $\chi u_t \in L^2(0,T;H_0^1(\omega_2))$, with
$\|\chi u_t\|_{L^2(0,T;H^1_0)}\le C(\chi)\big(a_*^{-1/2}+T^{1/2}\big)E_0^{1/2}$;
\item[\rm(ii)] $\partial_t(\chi u_t)\in L^2(0,T;H^{-1}(\omega_2))$, with a bound depending only on
$E_0$, $a_*$, $\chi$, $T$;
\item[\rm(iii)] $\chi u_t \in C\big([0,T];L^2(\omega_2)\big)$;
\item[\rm(iv)] $\chi u \in AC\big([0,T];H_0^1(\omega_2)\big)\subset C\big([0,T];H^1(\omega_2)\big)$.
\end{enumerate}
All four statements hold on the \emph{closed} interval $[0,T]$, with constants independent of
$T^*$.
\end{lemma}

\begin{proof}
\textbf{(i).} By \eqref{eq:astar}, $a\ge a_*$ on $\operatorname{supp}\chi\subset\omega_2$, so
\eqref{eq:two-bounds} gives
\begin{equation}\label{eq:grad-ut}
\int_0^{T}\!\!\int_{\omega_2}|\nabla u_t|^2\,dx\,dt
\;\le\; \frac1{a_*}\int_0^{T}\!\!\int_{\omega_2}a|\nabla u_t|^2\,dx\,dt \;\le\; \frac{E_0}{a_*}.
\end{equation}
Since $\nabla(\chi u_t)=\chi\nabla u_t+u_t\nabla\chi$ and
$\|u_t\|_{L^2(0,T;L^2)}\le T^{1/2}\sqrt{2E_0}$ by \eqref{eq:two-bounds},
\begin{eqnarray*}
&&\|\nabla(\chi u_t)\|_{L^2(0,T;L^2)}\le \|\nabla u_t\|_{L^2((0,T)\times\omega_2)}
+\|\nabla\chi\|_{L^\infty}\|u_t\|_{L^2(0,T;L^2)}\\
&&\le \Big(\frac{E_0}{a_*}\Big)^{1/2}+\|\nabla\chi\|_{L^\infty}T^{1/2}\sqrt{2E_0}.
\end{eqnarray*}
As $\chi$ has compact support in $\omega_2$, $\chi u_t(t)\in H_0^1(\omega_2)$ for a.e.\ $t$,
proving (i).

\textbf{(ii).} Since $\chi$ is independent of $t$, $\partial_t(\chi u_t)=\chi u_{tt}$, and from
\eqref{eq:main},
\begin{equation}\label{eq:three-terms}
\chi u_{tt} \;=\; \underbrace{\chi\Delta u}_{(\mathrm a)} \;+\; \underbrace{\chi\dive(a\nabla u_t)}_{(\mathrm b)} \;-\; \underbrace{\chi u^5}_{(\mathrm c)}
\qquad\text{in }\mathcal D'\big((0,T)\times\omega_2\big).
\end{equation}
We check each term lies in $L^2(0,T;H^{-1}(\omega_2))$.

\emph{(a)} Writing $\chi\Delta u=\Delta(\chi u)-2\nabla\chi\cdot\nabla u-(\Delta\chi)u$, we have
$\chi u\in L^\infty(0,T;H_0^1)$ so $\Delta(\chi u)\in L^\infty(0,T;H^{-1})$, while
$\nabla\chi\cdot\nabla u$ and $(\Delta\chi)u$ lie in $L^\infty(0,T;L^2)\hookrightarrow
L^\infty(0,T;H^{-1})$. All three are bounded by $C(\chi)\sqrt{2E_0}$.

\emph{(b)} Writing $\chi\dive(a\nabla u_t)=\dive\big(\chi a\nabla u_t\big)-a\,\nabla\chi\cdot\nabla u_t$,
the first term is the divergence of a field bounded in $L^2((0,T)\times\omega_2)$ by
$\|a\|_{L^\infty}(E_0/a_*)^{1/2}$ (using \eqref{eq:grad-ut}), hence lies in $L^2(0,T;H^{-1})$ with
the same bound; the second lies in $L^2((0,T)\times\omega_2)\hookrightarrow L^2(0,T;H^{-1})$,
bounded by $\|a\|_{L^\infty}\|\nabla\chi\|_{L^\infty}(E_0/a_*)^{1/2}$.

\emph{(c)} Since $H_0^1(\Om)\hookrightarrow L^6(\Om)$ in dimension three,
$u\in L^\infty(0,T;L^6)$ gives $u^5\in L^\infty(0,T;L^{6/5})$, and $L^{6/5}=(L^6)'\hookrightarrow
H^{-1}$; thus $\chi u^5\in L^\infty(0,T;H^{-1})\hookrightarrow L^2(0,T;H^{-1})$ for finite $T$,
with bound $C\,T^{1/2}E_0^{5/2}$.

Adding the three contributions proves (ii). \emph{Every bound used is an energy-class quantity;
no $H^2$ regularity of $u$, and no bound on $\nabla u_t$ outside $\omega_2$, has been invoked.}


\textbf{(iii).} Set $w:=\chi u_t$. By (i)--(ii), $w\in L^2(0,T;H_0^1(\omega_2))$ and \\
$w'\in L^2(0,T;H^{-1}(\omega_2))$. For the Gelfand triple
$H_0^1(\omega_2)\hookrightarrow L^2(\omega_2)\hookrightarrow H^{-1}(\omega_2)$, the
Lions--Magenes embedding theorem gives $w\in C([0,T];L^2(\omega_2))$ (after modification on a
null set), together with the estimate
$\|w\|_{C([0,T];L^2)}\le C\big(\|w\|_{L^2(H_0^1)}+\|w'\|_{L^2(H^{-1})}\big)$.


\textbf{(iv).} Since $\partial_t(\chi u)=\chi u_t$ and, by (i) and Cauchy--Schwarz,
\[
\int_0^{T}\big\|\chi u_t(t)\big\|_{H^1(\omega_2)}\,dt \;\le\; T^{1/2}\,\|\chi u_t\|_{L^2(0,T;H^1)} \;<\;\infty,
\]
we have $\chi u\in W^{1,1}(0,T;H_0^1(\omega_2))$, hence $\chi u$ is absolutely continuous with
values in $H_0^1(\omega_2)$ on $[0,T]$; in particular it is continuous there.
\end{proof}

\end{document}
